\printconcepts

\exercise{Use your own words to describe how sequences and series are related.}{Answers will vary.}

\exercise{Use your own words to define a \emph{partial sum}.}{Answers will vary.}

\exercise{Given a series $\ds \sum_{n=1}^\infty a_n$, describe the two sequences related to the series that are important.}{One sequence is the sequence of terms $\{a_n\}$. The other is the sequence of $n^{\text{th}}$ partial sums, $\{S_n\} = \{\sum_{i=1}^n a_i\}$.}

\exercise{Use your own words to explain what a geometric series is.}{Answers will vary.}

\exercise{T/F: If $\{a_n\}$ is convergent, then $\ds \sum_{n=1}^\infty a_n$ is also convergent.}{F}

\exercise{T/F: If $\{a_n\}$ converges to 0, then $\ds\sum_{n=0}^\infty a_n$ converges.}{F}

\printproblems

\input{exercises/08-02-exset-01}

%\input{exercises/08-02-exset-03}

\input{exercises/08-02-exset-04}

\exercisesetinstructions{, a series is given.
\begin{enumerate}
\item Find a formula for $S_n$, the $n^{\text{th}}$ partial sum of the series.
\item Determine whether the series converges or diverges. If it converges, state what it converges to.
\end{enumerate}}

\exercise{$\ds \sum_{n=0}^\infty \frac{1}{4^n}$}{\mbox{}\\[-2\baselineskip]\parbox[t]{\linewidth}{\begin{enumext}
\item	$S_n = \frac{1-(1/4)^n}{3/4}$
\item	Converges to $4/3$.
\end{enumext}}}

\exercise{$\ds 1^3+2^3+3^3+4^3+\dotsb$}{\mbox{}\\[-2\baselineskip]\parbox[t]{\linewidth}{\begin{enumext}
\item	$S_n = \left(\frac{n(n+1)}{2}\right)^2$
\item	Diverges
\end{enumext}}}

\exercise{$\ds \sum_{n=1}^\infty (-1)^n n$}{\mbox{}\\[-2\baselineskip]\parbox[t]{\linewidth}{\begin{enumext}
\item	$S_n = \begin{cases}-\frac{n+1}{2} & $n$ \text{ is odd }\\
	\frac{n}{2} & $n$ \text{ is even}
	\end{cases}$
\item	Diverges
\end{enumext}}}

\exercise{$\ds \sum_{n=0}^\infty \frac{5}{2^n}$}{\mbox{}\\[-2\baselineskip]\parbox[t]{\linewidth}{\begin{enumext}
\item	$S_n = 5\frac{1-1/2^n}{1/2}$
\item	Converges to 10.
\end{enumext}}}

\exercise{$\ds \sum_{n=1}^\infty e^{-n}$}{\mbox{}\\[-2\baselineskip]\parbox[t]{\linewidth}{\begin{enumext}
\item		$S_n = \frac{1-(1/e)^{n+1}}{1-1/e}$.
\item		Converges to $1/(1-1/e) = e/(e-1)$.
\end{enumext}}}

\exercise{$\ds 1-\frac13+\frac19-\frac{1}{27}+\frac{1}{81}+\dotsb$}{\mbox{}\\[-2\baselineskip]\parbox[t]{\linewidth}{\begin{enumext}
\item	$S_n = \frac{1-(-1/3)^n}{4/3}$
\item	Converges to $3/4$.
\end{enumext}}}

\exercise{$\ds \sum_{n=1}^\infty \frac{1}{n(n+1)}$}{\mbox{}\\[-2\baselineskip]\parbox[t]{\linewidth}{\begin{enumext}
\item		With partial fractions, $a_n = \frac1n-\frac1{n+1}$.\\
Thus $S_n = 1-\frac1{n+1}$.
\item		Converges to 1.
\end{enumext}}}

\exercise{$\ds \sum_{n=1}^\infty \frac{3}{n(n+2)}$}{\mbox{}\\[-2\baselineskip]\parbox[t]{\linewidth}{\begin{enumext}
\item		With partial fractions, $a_n = \frac32\left(\frac1n-\frac1{n+2}\right)$.\\
Thus $S_n = \frac32\left(\frac32-\frac1{n+1}-\frac1{n+2}\right)$.
\item		Converges to 9/4
\end{enumext}}}

\exercise{$\ds \sum_{n=1}^\infty \frac{1}{(2n-1)(2n+1)}$}{\mbox{}\\[-2\baselineskip]\parbox[t]{\linewidth}{\begin{enumext}
\item	Use partial fraction decomposition to recognize the telescoping series: $a_n = \frac12\left(\frac{1}{2n-1}-\frac{1}{2n+1}\right)$, so that $S_n=\frac12\left(1-\frac1{2n+1}\right) = \frac{n}{2n+1}$.
\item	Converges to $1/2$.
\end{enumext}}}

\exercise{$\ds \sum_{n=1}^\infty \ln \left(\frac{n}{n+1}\right)$}{\mbox{}\\[-2\baselineskip]\parbox[t]{\linewidth}{\begin{enumext}
\item	$S_n = \ln \big(1/(n+1)\big)$
\item	Diverges (to $-\infty$).
\end{enumext}}}

\exercise{$\ds \sum_{n=1}^\infty \frac{2n+1}{n^2(n+1)^2}$}{\mbox{}\\[-2\baselineskip]\parbox[t]{\linewidth}{\begin{enumext}
\item	$S_n = 1-\frac{1}{(n+1)^2}$
\item	Converges to 1.
\end{enumext}}}

\exercise{$\ds \frac{1}{1\cdot 4}+\frac1{2\cdot5}+\frac1{3\cdot6}+\frac1{4\cdot7}+\dotsb$}{\mbox{}\\[-2\baselineskip]\parbox[t]{\linewidth}{\begin{enumext}
\item	$a_n = \frac1{n(n+3)}$; using partial fractions,\\
the resulting telescoping sum reduces to\\
$S_n = \frac13\left(1+\frac12+\frac13-\frac1{n+1}-\frac1{n+2}-\frac1{n+3}\right)$
\item	Converges to $11/18$.
\end{enumext}}}

\exercise{$\ds 2+\left(\frac12+\frac13\right) + \left(\frac14+\frac19\right) + \left(\frac18+\frac1{27}\right)+\dotsb$}{\mbox{}\\[-2\baselineskip]\parbox[t]{\linewidth}{\begin{enumext}
\item	$a_n = 1/2^n+1/3^n$ for $n\geq 0$. Thus $S_n = \frac{1-1/2^2}{1/2}+\frac{1-1/3^n}{2/3}$.
\item		Converges to $2+3/2 = 7/2$.
\end{enumext}}}

\exercise{$\ds \sum_{n=2}^\infty \frac{1}{n^2-1}$}{\mbox{}\\[-2\baselineskip]\parbox[t]{\linewidth}{\begin{enumext}
\item		With partial fractions, $a_n = \frac12\left(\frac1{n-1}-\frac1{n+1}\right)$. Thus $S_n = \frac12\left(3/2-\frac1n-\frac{1}{n+1}\right)$.
\item		Converges to 3/4.
\end{enumext}}}

\exercise{$\ds \sum_{n=0}^\infty \big(\sin 1\big)^n$}{\mbox{}\\[-2\baselineskip]\parbox[t]{\linewidth}{\begin{enumext}
\item		$S_n=\frac{1-(\sin 1)^{n+1}}{1-\sin 1}$
\item		Converges to $\frac{1}{1-\sin 1}$.
\end{enumext}}}

\exercise{$\ds \sum_{n=1}^\infty \biggl(\frac{2}{n(n+2)} +\frac{5}{4^n}\biggr)$}{\mbox{}\\[-2\baselineskip]\parbox[t]{\linewidth}{\begin{enumext}
\item		$S_n=1+\frac12-\frac{1}{n+1}-\frac{1}{n+2}+\frac{\frac54(1-(\frac14)^n)}{1-\frac14}$
\item		Converges to $\frac{19}{6}$.
\end{enumext}}}

\exercisesetend


\input{exercises/08-02-exset-05}

%\exercise{\label{08_02_ex_24} Break the Harmonic Series into the sum of the odd and even terms:
%\[\sum_{n=1}^\infty \frac1n = \sum_{n=1}^\infty \frac{1}{2n-1}+\sum_{n=1}^\infty \frac{1}{2n}.\] The goal is to show that each of the series on the right diverge.
%		\begin{enumext}
%		\item		Show why $\ds \sum_{n=1}^\infty \frac{1}{2n-1}>\sum_{n=1}^\infty \frac{1}{2n}$. 
%		
%		(Compare each $n^{\text{th}}$ partial sum.)
%		\item		Show why $\ds\sum_{n=1}^\infty \frac{1}{2n-1}<1+\sum_{n=1}^\infty \frac{1}{2n}$
%		\item		Explain why (a) and (b) demonstrate that the series of odd terms is convergent, if, and only if, the series of even terms is also convergent. (That is, show both converge or both diverge.)
%		\item		Explain why knowing the Harmonic Series is divergent determines that the even and odd series are also divergent.
%		\end{enumext}
%}{\begin{enumext}
%\item		The $n^{\text{th}}$ partial sum of the odd series is $1+\frac13+\frac15+\dotsb+\frac{1}{2n-1}$. The $n^{\text{th}}$ partial sum of the even series is $\frac12+\frac14 + \frac16 + \dotsb +\frac1{2n}$. Each term of the even series is less than the corresponding term of the odd series, giving us our result.
%\item		The $n^{\text{th}}$ partial sum of the odd series is $1+\frac13+\frac15+\dotsb+\frac1{2n-1}$. The $n^{\text{th}}$ partial sum of 1 plus the even series is $1+\frac12+\frac14+\dotsb + \frac{1}{2(n-1)}$. Each term of the even series is now greater than or equal to the corresponding term of the odd series, with equality only on the first term. This gives us the result.
%\item		If the odd series converges, the work done in (a) shows the even series converges also. (The sequence of the $n^{\text{th}}$ partial sum of the even series is bounded and monotonically increasing.) Likewise, (b) shows that if the even series converges, the odd series will, too. Thus if either series converges, the other does. 
%
%Similarly, (a) and (b) can be used to show that if either series diverges, the other does, too.
%\item		If both the even and odd series converge, then their sum would be a convergent series. This would imply that the Harmonic Series, their sum, is convergent. It is not. Hence each series diverges.
%\end{enumext}
%}

\exercisesetinstructions{, use \autoref{thm:series_nth_term} to show the given series diverges.}

\exercise{$\ds \sum_{n=1}^\infty \frac{3n^2}{n(n+2)}$}{$\ds \lim_{n\to\infty}a_n = 3$; by \autoref{thm:series_nth_term} the series diverges.}

\exercise{$\ds \sum_{n=1}^\infty \frac{2^n}{n^2}$}{$\ds \lim_{n\to\infty}a_n = \infty$; by \autoref{thm:series_nth_term} the series diverges.}

\exercise{$\ds \sum_{n=1}^\infty \frac{n!}{10^n}$}{$\ds \lim_{n\to\infty}a_n = \infty$; by \autoref{thm:series_nth_term} the series diverges.}

\exercise{$\ds \sum_{n=1}^\infty \frac{5^n-n^5}{5^n+n^5}$}{$\ds \lim_{n\to\infty}a_n = 1$; by \autoref{thm:series_nth_term} the series diverges.}

\exercise{$\ds \sum_{n=1}^\infty \frac{2^n+1}{2^{n+1}}$}{$\ds \lim_{n\to\infty}a_n = 1/2$; by \autoref{thm:series_nth_term} the series diverges.}

\exercise{$\ds \sum_{n=1}^\infty \left(1+\frac1n\right)^n$}{$\ds \lim_{n\to\infty}a_n = e$; by \autoref{thm:series_nth_term} the series diverges.}

\exercisesetend

\exercise{Show the series $\ds \sum_{n=1}^\infty \frac{n}{(2n-1)(2n+1)}$ diverges.}{Using partial fractions, we can show that $a_n =
\frac14\left(\frac1{2n-1}+\frac{1}{2n+1}\right)$. The series is effectively twice the sum of the odd terms of the Harmonic Series which was shown to diverge in \autoref{ex_harm_series}. Thus this series diverges.}

\exercise{Rewrite $0.12 12 \overline{12} \dots$ as an infinite series and then express the sum as the quotient of two integers (thus showing that it is a rational number --- this method can be generalized to show that every repeating decimal is rational).}{$\sum_{n=1}^\infty12(100)^{-n}=\frac{12}{100}\cdot\frac1{1-\frac1{100}}=\frac{12}{100-1}=\frac{12}{99}=\frac4{33}$}

\exercise{A ball falling from a height of $h$ meters is known to rebound to a height of $rh$ meters, where the proportionality constant $r$ satisfies $0<r<1$. Find the total distance traveled (vertically) by the ball if it is dropped initially from a height of 2 meters.}{$2\left(\frac{1+r}{1-r}\right)$}
